\section{Introduction} \label{sec:intro}

Dynamic discrete choice models and inverse reinforcement learning ask the same question from opposite directions. The structural econometrics tradition starts from a known reward function with unknown parameters and recovers those parameters by maximum likelihood. The inverse reinforcement learning tradition starts from an unknown reward and recovers it from observed expert behavior. Both rely on the same underlying primitive, a Markov decision process with a forward-looking decision maker. Both target the same downstream goal, namely to forecast behavior under counterfactual environments. Yet the two literatures have grown up in different communities with different software stacks. \citet{rust1987} introduced the nested fixed point algorithm in \proglang{Fortran}. \citet{hotz1993} developed the conditional choice probability estimator that has since been re-implemented in \proglang{Stata}, \proglang{Matlab}, and \proglang{R}. \citet{ziebart2008} introduced maximum entropy inverse reinforcement learning in \proglang{Python} with bespoke optimization code. \citet{fu2018} introduced adversarial inverse reinforcement learning on top of \proglang{TensorFlow}. \citet{garg2021} introduced inverse soft Q-learning on top of \proglang{PyTorch}. The result is a fragmented ecosystem in which methods that solve nearly the same problem cannot exchange data, do not share an inference layer, and resist apples-to-apples comparison.

The cost of this fragmentation falls on practitioners. An applied economist who fits a structural model in \proglang{Stata} cannot easily test whether a deep inverse reinforcement learning method would recover a richer reward on the same panel. A reinforcement learning researcher who trains an adversarial reward learner in \proglang{PyTorch} has no path to standard errors, identification diagnostics, or counterfactual welfare elasticity. A graduate student who wants to compare twelve estimators on the same dataset must reimplement most of them. The estimators themselves are well documented in their respective papers. What is missing is the shared infrastructure that turns each estimator into a usable tool for inference and policy analysis.

\pkg{econirl} is an open-source \proglang{Python} library that addresses this fragmentation. It implements twelve estimators spanning structural maximum likelihood, conditional choice probability, sieve and neural value approximation, maximum entropy inverse reinforcement learning, distribution matching, inverse soft Q-learning, adversarial inverse reinforcement learning, and behavioral cloning. Every estimator implements the same \code{Estimator} protocol and returns the same \code{EstimationSummary} object. This single contract enables a shared post-estimation pipeline. Standard errors are available for any estimator that exposes a Hessian, by four interchangeable methods including asymptotic, robust sandwich, nonparametric bootstrap, and clustered. Identification diagnostics are computed from the same Hessian. Counterfactual simulation, profile likelihood, hypothesis testing, and visualization all consume the unified result object. The user picks an estimator and the rest of the pipeline carries through unchanged.

The library is built on \pkg{JAX} \citep{jax2018github}. Automatic differentiation produces the gradients and Hessians that the inference layer needs without hand-coded analytic derivatives. Just-in-time compilation accelerates the inner loops of the structural estimators. Vectorized mapping over individual trajectories enables bootstrap resampling at small marginal cost. Graphics processing unit support is available for the neural estimators on larger state spaces. The library is distributed under the MIT license and works on any machine with \proglang{Python} 3.10 or newer.

Listing~\ref{lst:teaser} demonstrates the user-facing workflow. The same eight lines that fit a nested fixed point estimator on the canonical Rust bus engine replacement panel would fit any of the other eleven estimators by substituting the class name. The returned \code{EstimationSummary} object carries point estimates, standard errors, identification diagnostics, the value function, the policy, and a \code{summary()} method that prints a regression-style table.

\begin{lstlisting}[caption={Fitting NFXP on the Rust bus panel and printing the result.},label={lst:teaser},language=Python,frame=single]
from econirl.datasets import load_rust_bus
from econirl.estimation import NFXP

panel = load_rust_bus()
result = NFXP(discount_factor=0.9999).estimate(panel)
print(result.summary())
print(result.identification.condition_number)
print(result.confidence_interval(level=0.95))
\end{lstlisting}

\subsection{Related software}

The \pkg{imitation} library \citep{gleave2022imitation} from the Center for Human-Compatible Artificial Intelligence is the most mature open-source package for inverse reinforcement learning. It implements behavioral cloning, generative adversarial imitation learning, adversarial inverse reinforcement learning, and a tabular variant of maximum causal entropy inverse reinforcement learning. The library is excellent within its scope but does not implement structural likelihood estimators, does not compute standard errors, does not run identification diagnostics, and does not include a counterfactual analysis pipeline. Its design assumes the \pkg{Stable Baselines 3} ecosystem and the \pkg{Gymnasium} environment interface. \pkg{econirl} borrows the direct linear solve for occupancy measures from \pkg{imitation} and the tabular maximum causal entropy framework but extends both to the infinite-horizon action-dependent feature setting that structural econometrics requires.

The \pkg{pyblp} package \citep{conlon2020pyblp} is the leading open-source tool for random coefficients logit demand estimation. It is mature and battle-tested but addresses a static demand setting rather than a dynamic decision problem. The R packages \pkg{mlogit} and \pkg{Rchoice} fit static discrete choice models with no inter-temporal substitution. None of these libraries cover dynamic decision making with reward learning.

The closest existing tool to \pkg{econirl} in the structural econometrics tradition is \pkg{respy} \citep{respy2020}, a \proglang{Python} library for structural estimation of Eckstein-Keane-Wolpin models. \pkg{respy} is specialized for finite-horizon labor supply problems and does not implement inverse reinforcement learning estimators. \pkg{econirl} aims to complement rather than replace these specialized tools by providing a single home for the broader family of dynamic discrete choice estimators and the inverse reinforcement learning estimators that share the underlying Markov decision process structure.

\subsection{Contributions}

The contribution of this paper is the design of a unified result object and a shared post-estimation pipeline that together turn twelve heterogeneous estimators into interchangeable parts of a single workflow. Section~\ref{sec:methods} lays out the unified notation, the families of estimators that the library covers, and the structural assumptions each family makes. Section~\ref{sec:design} describes the \code{Estimator} protocol, the \code{EstimationSummary} contract, the inference layer, the counterfactual layer, and the JAX backbone. Section~\ref{sec:illustrations} works through three end-to-end examples on real data. The first fits a nested fixed point estimator on the Rust bus engine replacement panel and walks through the full inference and counterfactual workflow. The second fits a maximum causal entropy inverse reinforcement learning estimator on the same panel and demonstrates that the unified pipeline supports both structural utility recovery and inverse reward recovery without code changes outside the estimator itself. The third fits a neural value approximation estimator on the larger Keane-Wolpin occupational choice panel and reports graphics processing unit timings. Section~\ref{sec:illustrations} closes with a cross-estimator benchmark table covering all twelve estimators. Section~\ref{sec:computational} reports the package versions and hardware used for the illustrations. Section~\ref{sec:summary} closes with a roadmap for continuous-action and dynamic-game extensions. Appendix~\ref{app:reproducing} maps every numbered figure and table in the paper to a runnable script.
